\section{Рекурсивная модель и переиспользование малых языковых моделей}\label{sec:slm_reuse}

Авторская рекурсивная схема обсуждается в литературе как tiny recursive model (TRM)~\cite{trm}. Последовательность внутренних состояний, распределений \(Q_{t,k}\) и голов времени описана в разд.~\ref{sec:recursive_process}; параметризация блоков \(E_{\phi}\), оператора \(U_\theta\) (в виде остаточного ядра \(\mathcal U_{\theta}\)) и возможная композиция с малой языковой моделью задаются формально в разд.~\ref{sec:model_architecture}. Ниже сосредоточены сравнение с методом LoRA и кратко зафиксировано переиспользование весов малой предобученной модели.

Смысл рекурсии состоит не в механическом увеличении числа слоёв поверх уже обученной сетки. Один и тот же оператор уточнения применяется к состоянию несколько раз, а число применений определяется данными через правило остановки над \(\widehat p_{t,k}\). Поэтому модель может тратить меньше вычислений на простые позиции и больше --- на сложные.

\subsection{Сравнение с LoRA}

Метод LoRA~\cite{hu2021} решает типичную инженерную задачу: не переписывая большие тензоры предобученной модели, повысить её качество на целевой задаче. Базовые веса остаются фиксированными, а в выбранных линейных слоях появляется учимое низкоранговое приращение вида \(\alpha\,B^{\top}A\). После обучения граф вывода остаётся обычным для трансформера: одна и та же фиксированная глубина сетки, одинаковая для всех позиций и всех предложений; меняется только сама функция, реализуемая сеткой, за счёт добавленных параметров.

Предлагаемая рекурсивная схема физически устроена иначе. Качество на позиции повышается не (или не только) за счёт перенастройки матриц внутри слоёв базы, а за счёт \emph{политики внутреннего времени}: одно и то же ядро \(\mathcal U_{\theta}\) может применяться разное число раз, а решение о продолжении или остановке принимается по головам \(\widehat p_{t,k}\) и правилу остановки (разд.~\ref{sec:recursive_process}). При замороженной малой языковой модели это отдельный вычислительный контур над её представлением --- проекция и рекурсия из разд.~\ref{sec:model_architecture} --- а не изменение топологии «одного прохода по трансформеру».

С практической стороны и LoRA, и рекурсия на замороженном чекпоинте отвечают на один и тот же вопрос: как выжать качество из \emph{той же} малой предобученной модели, доплатив сравнимо умеренным числом обучаемых параметров или модулей. Без режима «только LoRA» контраст «рекурсивная надстройка» против «замороженная база без малоранговой адаптации трансформера» остаётся неоднозначным: вклад малоранговой параметрической подгонки к данным смешался бы с вкладом переменной глубины цикла. Добавление LoRA в решётку не означает, что методы конкурируют концептуально; это явный второй базовый способ улучшить данную модель, относительно которого можно оценить, оправдана ли именно управляемая глубина.


\subsection{Переиспользование весов готовой модели}

Предлагается использовать TRM --- проектор \(P_{\psi}\), ядро \(\mathcal U_{\theta}\), прогностическая и временная головы (разд.~\ref{sec:model_architecture}) --- как надстройку над полностью замороженными весами малой предобученной модели Falcon-H1-Tiny порядка \(90~\mathrm{M}\) параметров (\texttt{tiiuae/Falcon-H1-Tiny-90M-Base}), получая признаки с её \texttt{last\_hidden\_state}. На том же чекпоинте для сравнения по качеству отдельно прогоняют вариант с LoRA~\cite{hu2021} без рекурсивного контура; объединять LoRA и рекурсию в одном прогоне в данной выпускной работе не предполагается.
